\chapter{第12章\texorpdfstring{\hspace{0.5em}}{ }波动学基础}

\section{机械波的产生和传播}
    \subsection{机械波的形成}
    机械振动在弹性介质中的传播过程称为机械波。要形成机械波，必须有做机械振动的波源和能够传播振动的弹性介质。波源振动时，由于介质各质点之间存在弹性联系，波源处的质点带动相邻质点振动，相邻质点又带动更远的质点，于是振动状态（相位）由近及远地传播出去，形成波动。

    需要强调：波动传播的是振动状态和能量，介质中各质点只在各自的平衡位置附近振动，并不随波迁移。所以"波是振动状态的传播，不是介质的迁移"。

    按质点振动方向与波传播方向的关系，机械波分为两类：
    \begin{enumerate}
        \item 横波：质点振动方向与波的传播方向垂直，如绳上的波。横波传播时介质发生切变，因此只能在固体中传播；
        \item 纵波：质点振动方向与波的传播方向平行，如声波。纵波传播时介质发生体变（压缩与拉伸），可在固体、液体和气体中传播。
    \end{enumerate}
    为描述波在空间的传播，引入波线、波面和波前：波线（射线）表示波的传播方向；波面上各点相位相同；波前是最前面的那个波面。在各向同性介质中，波线与波面垂直。波面为平面的波称平面波，波面为球面的波称球面波。

    \subsection{描述波动的物理量}
    描述波动常用以下物理量：
    \begin{enumerate}
        \item 波长 $\lambda$：同一波线上相位差为 $2\pi$ 的两个相邻质点间的距离，即一个周期内波传播的距离；
        \item 周期 $T$：波传播一个波长所需的时间，也等于介质质点振动的周期，由波源决定；
        \item 频率 $\nu=\dfrac{1}{T}$：单位时间内波向前传播的完整波数，由波源决定；
        \item 波速 $u$：振动状态（相位）在介质中传播的速度，由介质的性质决定。
    \end{enumerate}
    三者之间的关系为
    \keyformula{u=\frac{\lambda}{T}=\lambda\nu}
    这是波动学的基本关系式。

    必须注意：波从一种介质进入另一种介质时，频率（周期）由波源决定而保持不变，波速由介质决定而改变，因此波长也要改变（$\lambda=u/\nu$）。这一点与"光的颜色由频率决定"是同一个道理。

    \corepoint{（1）波速是振动状态的传播速度，不是介质质点的振动速度，二者不可混淆；（2）频率由波源决定，波速由介质决定，波长随二者而变；（3）横波只能在固体中传播，纵波在固、液、气中都能传播，因此声波（纵波）能在空气中传播而绳波（横波）不能。}

\section{平面简谐波的波函数}
    \subsection{波函数的建立}
    波源作简谐振动、且在介质中传播时波形不改变的波称简谐波。设一平面简谐波以波速 $u$ 沿 $x$ 轴正方向传播，原点 $O$ 处质点的振动表达式为
    \[
    y(0,t)=A\cos(\omega t+\varphi),
    \]
    则 $x$ 处质点的振动比原点处落后一段时间 $\Delta t=x/u$（因为振动状态从原点传到 $x$ 处需要这段时间），故 $x$ 处质点在 $t$ 时刻的位移等于原点处质点在 $t-\Delta t$ 时刻的位移：
    \[
    y(x,t)=A\cos\left[\omega\left(t-\frac{x}{u}\right)+\varphi\right].
    \]
    利用 $u=\lambda\nu$、$\omega=2\pi\nu$，可写成几种等价形式：
    \keyformula{y(x,t)=A\cos\left[\omega\left(t-\frac{x}{u}\right)+\varphi\right]=A\cos\left[2\pi\left(\frac{t}{T}-\frac{x}{\lambda}\right)+\varphi\right]=A\cos(\omega t-kx+\varphi)}
    其中
    \[
    k=\frac{2\pi}{\lambda}
    \]
    称角波数（波数）。这三式都是沿 $x$ 轴正方向传播的平面简谐波的波函数。

    若波沿 $x$ 轴负方向传播，则 $x$ 处质点的振动超前原点，波函数中 $x/u$ 前取正号：
    \[
    y(x,t)=A\cos\left[\omega\left(t+\frac{x}{u}\right)+\varphi\right].
    \]
    若已知波线上某点 $x_0$ 的振动表达式为 $y=A\cos(\omega t+\varphi_0)$，则波函数为
    \[
    y(x,t)=A\cos\left[\omega\left(t-\frac{x-x_0}{u}\right)+\varphi_0\right],
    \]
    即把"某点的相位"沿传播方向逐点推迟。

    \subsection{波函数的物理意义}
    波函数 $y(x,t)$ 含有两个自变量，从不同角度理解它：
    \begin{enumerate}
        \item 固定 $x=x_1$：$y=A\cos(\omega t-kx_1+\varphi)$ 表示该处质点作简谐振动，其初相为 $\varphi-kx_1$，说明不同位置的质点振动相位不同，$x$ 越大相位越落后；
        \item 固定 $t=t_1$：$y=A\cos(kx-\omega t_1-\varphi)$ 表示该时刻各质点的位移分布，即波形曲线（余弦曲线）；
        \item $t$ 与 $x$ 都变：波形随时间沿传播方向平移，在 $\Delta t$ 时间内波形平移 $\Delta x=u\Delta t$，这说明波是振动状态以速度 $u$ 向前传播。
    \end{enumerate}
    同一时刻波线上相距 $\Delta x$ 的两点，其相位差为
    \[
    \Delta\varphi=-k\Delta x=-\frac{2\pi}{\lambda}\Delta x,
    \]
    负号表示沿传播方向相位依次落后。相距一个波长时相位差为 $2\pi$，相距半个波长时相位差为 $\pi$（反相）。

    \corepoint{（1）波函数中 $x/u$ 前的符号决定传播方向：负号表示沿 $+x$ 传播，正号表示沿 $-x$ 传播；（2）"落后"是指相位落后，因为振动从波源传来需要时间；（3）写波函数时必须先确定原点（或已知点）的振动表达式和波的传播方向，二者缺一不可。}

\section{波动方程与波速}
    \subsection{物体的弹性形变}
    机械波在介质中传播依赖于介质的弹性。常见的三种基本形变为：
    \begin{enumerate}
        \item 线变（拉伸或压缩）：长 $l$、截面积 $S$ 的棒受拉力 $F$ 伸长 $\Delta l$，正应力 $F/S$ 与线应变 $\Delta l/l$ 之比为杨氏模量
        \[
        E=\frac{F/S}{\Delta l/l};
        \]
        \item 切变：立方体受切向力作用，切应力与切应变之比为切变模量 $G$；
        \item 体变：体积为 $V$ 的介质受压，压强增量 $\Delta p$ 与体积相对变化 $-\Delta V/V$ 之比为体积模量
        \[
        K=-\frac{\Delta p}{\Delta V/V}.
        \]
    \end{enumerate}

    \subsection{波动方程}
    将平面简谐波的波函数 $y=A\cos[\omega(t-x/u)+\varphi]$ 分别对 $x$ 和 $t$ 求二阶偏导，可得
    \[
    \frac{\partial^{2}y}{\partial x^{2}}=\frac{1}{u^{2}}\frac{\partial^{2}y}{\partial t^{2}},
    \]
    称为一维波动方程（平面波波动方程）。它是一切以速度 $u$ 沿 $x$ 方向传播的波（不论机械波还是电磁波）所满足的普遍方程，平面简谐波的波函数只是它的一个特解。

    \subsection{决定波速的因素}
    由波动方程可以证明，波速由介质的弹性模量和密度决定，一般可写成
    \[
    u=\sqrt{\frac{\text{弹性模量}}{\text{密度}}}.
    \]
    具体地：
    \begin{table}[htbp]
    \centering
    \caption{几种介质中的波速}
    \begin{tabular}{ll}
    \toprule
    波的类型与介质 & 波速 \\
    \midrule
    绳（弦）上的横波 & $u=\sqrt{\dfrac{F}{\rho_l}}$（$F$ 为张力，$\rho_l$ 为线密度） \\
    固体中的纵波 & $u=\sqrt{\dfrac{E}{\rho}}$（$E$ 为杨氏模量） \\
    固体中的横波 & $u=\sqrt{\dfrac{G}{\rho}}$（$G$ 为切变模量） \\
    液体、气体中的纵波 & $u=\sqrt{\dfrac{K}{\rho}}$（$K$ 为体积模量） \\
    理想气体中的声波 & $u=\sqrt{\dfrac{\gamma RT}{M}}$（$\gamma$ 为比热比，$M$ 为摩尔质量） \\
    \bottomrule
    \end{tabular}
    \end{table}

    可见，波速只由介质的性质（弹性模量、密度、温度）决定，与波的频率、振幅无关。同一介质中，纵波波速一般大于横波波速（因为 $E>G$），这解释了地震时先感到上下颠动（纵波）后感到左右摇晃（横波）的现象。

    \corepoint{波速由介质决定、频率由波源决定，二者不能混为一谈。波从一种介质进入另一种介质时 $u$ 变、$\lambda$ 变而 $\nu$ 不变。}

\section{波的能量}
    \subsection{平面简谐纵波传播时介质元的能量}
    在介质中取一体积元 $\mathrm{d}V$、质量 $\mathrm{d}m=\rho\mathrm{d}V$，当波传到该处时，其振动速度为
    \[
    v=\frac{\partial y}{\partial t}=-A\omega\sin\left[\omega\left(t-\frac{x}{u}\right)+\varphi\right].
    \]
    该体积元的振动动能为
    \[
    \mathrm{d}E_\mathrm{k}=\frac{1}{2}\mathrm{d}m\,v^{2}
    =\frac{1}{2}\rho\mathrm{d}V\,A^{2}\omega^{2}\sin^{2}\left[\omega\left(t-\frac{x}{u}\right)+\varphi\right].
    \]
    由于体积元同时发生形变，还具有弹性势能，可以证明其弹性势能与动能相等且同步变化：
    \[
    \mathrm{d}E_\mathrm{p}=\mathrm{d}E_\mathrm{k}.
    \]
    因此该体积元的总机械能为
    \[
    \mathrm{d}E=\mathrm{d}E_\mathrm{k}+\mathrm{d}E_\mathrm{p}
    =\rho\mathrm{d}V\,A^{2}\omega^{2}\sin^{2}\left[\omega\left(t-\frac{x}{u}\right)+\varphi\right].
    \]

    这与单个简谐振子的能量有本质区别：单个振子的动能与势能相互转化、此消彼长，总能量守恒；而介质元的动能与势能同相变化，同时达到最大、同时为零，总能量不守恒，它随时间作周期性变化，能量不断从波源向前传递，这正是波传播能量的表现。

    \subsection{波的能量密度和能流密度}
    单位体积介质中的波的能量称为能量密度：
    \[
    w=\frac{\mathrm{d}E}{\mathrm{d}V}=\rho A^{2}\omega^{2}\sin^{2}\left[\omega\left(t-\frac{x}{u}\right)+\varphi\right].
    \]
    它随时间变化，其一个周期内的平均值为
    \[
    \bar{w}=\frac{1}{2}\rho A^{2}\omega^{2}.
    \]

    单位时间内通过某一面积的能量称能流（功率）$P=wuS$（$S$ 为垂直于传播方向的面积）。单位时间内通过垂直于传播方向的单位面积的能量称能流密度，也称波的强度（波强）：
    \keyformula{I=\frac{P}{S}=\bar{w}u=\frac{1}{2}\rho A^{2}\omega^{2}u}
    单位是 $\mathrm{W/m^{2}}$。能流密度是描述波的强弱的重要物理量，它与振幅的平方、频率的平方及波速成正比。

    \subsection{波的吸收}
    波在介质中传播时，由于介质的吸收和散射，波的强度随传播距离逐渐减小。实验表明，强度按指数规律衰减：
    \[
    I=I_0\mathrm{e}^{-\alpha x},
    \]
    其中 $\alpha$ 称吸收系数，与介质的性质和波的频率有关。频率越高，吸收一般越强，这就是超声波在空气中传播不远、而低频声波能传播很远的原因。

    \subsection{球面波}
    球面波从点波源向各方向均匀传播，同一波面上的强度相同。由能量守恒，通过半径为 $r_1$、$r_2$ 的两个波面的总功率相等，即 $I_1\cdot4\pi r_1^{2}=I_2\cdot4\pi r_2^{2}$，故
    \[
    I\propto\frac{1}{r^{2}}.
    \]
    又因 $I\propto A^{2}$，得振幅与距离成反比：
    \[
    A_1r_1=A_2r_2,
    \]
    即球面波的振幅随距离反比衰减。平面波则因波面面积不变，振幅不随距离改变。

    \corepoint{（1）波的能量密度中动能与势能同相变化，这是波与单个振子的根本区别；（2）能流密度 $I=\frac{1}{2}\rho A^{2}\omega^{2}u$ 中三个量都要用国际单位；（3）球面波振幅按 $1/r$ 衰减、强度按 $1/r^{2}$ 衰减，二者不要混淆。}

\section{惠更斯原理}
    \subsection{惠更斯原理}
    惠更斯于 1690 年提出：介质中波前上的每一点都可以看作发射子波的新波源，这些子波在介质中向前传播，在其后某一时刻，这些子波的包迹（包络面）就是该时刻新的波前。这就是惠更斯原理。

    惠更斯原理的实质是：波前上各点都可作为新的波源，波的传播方向就是子波包迹的法线方向。它对任何波动（机械波、电磁波）都适用，只要能确定某一时刻的波前，就能确定以后各时刻的波前。

    \subsection{惠更斯原理的应用}
    \begin{enumerate}
        \item 解释波的衍射：波遇到障碍物或狭缝时，若缝宽与波长可比，缝上各点作为子波源发出的子波会向障碍物后面传播，使波偏离直线传播而绕到障碍物后面，形成衍射；
        \item 解释波的反射和折射：由子波包迹可导出反射定律（入射角等于反射角）和折射定律（$\dfrac{\sin i}{\sin r}=\dfrac{u_1}{u_2}$）；
        \item 说明波在两种介质界面上传播方向改变的原因：波速改变导致波前方向改变。
    \end{enumerate}

    惠更斯原理的局限在于它只能说明波的传播方向（几何性质），不能说明波的强度分布，也不能解释干涉、衍射条纹的明暗分布，这些问题需要由惠更斯-菲涅耳原理（考虑子波的相干叠加）来解决。

    \corepoint{惠更斯原理中的"子波源"是波前上的点，不是波源本身；该原理只能确定波前的几何形状和传播方向，不能给出振幅和强度信息。}

\section{波的叠加原理和波的干涉}
    \subsection{波的叠加原理}
    实验表明：几列波在同一介质中传播而相遇时，在相遇区域内任一点处，介质质点的位移等于各列波单独存在时在该点所引起的位移的矢量和；相遇后各列波仍按各自原来的方向、频率和波形继续传播，互不影响。这称为波的叠加原理，也称波的独立传播原理。

    叠加原理成立的条件是介质的形变很小、满足线性关系。当波的振幅很大（如爆炸产生的冲击波）时，介质表现出非线性，叠加原理不再成立。

    \subsection{波的干涉}
    两列波在空间相遇，若在相遇区域内某些点的振动始终加强、某些点的振动始终减弱，形成稳定的振幅分布，这种现象称为波的干涉。能产生干涉现象的两列波称为相干波，其波源称相干波源。相干条件为：
    \begin{enumerate}
        \item 频率相同；
        \item 振动方向相同；
        \item 相位差恒定。
    \end{enumerate}

    设两相干波源 $S_1$、$S_2$ 的振动分别为 $y_1=A_1\cos(\omega t+\varphi_1)$、$y_2=A_2\cos(\omega t+\varphi_2)$，它们到空间某点 $P$ 的距离分别为 $r_1$、$r_2$，则两波在 $P$ 点的相位差为
    \[
    \Delta\varphi=(\varphi_2-\varphi_1)-\frac{2\pi}{\lambda}(r_2-r_1).
    \]
    两波在 $P$ 点叠加后的合振幅为
    \keyformula{A=\sqrt{A_1^{2}+A_2^{2}+2A_1A_2\cos\Delta\varphi}}
    由此得干涉加强和减弱的条件：
    \begin{enumerate}
        \item 当 $\Delta\varphi=\pm2k\pi$（$k=0,1,2,\cdots$）时，$A=A_1+A_2$，干涉加强（相长）；
        \item 当 $\Delta\varphi=\pm(2k+1)\pi$ 时，$A=|A_1-A_2|$，干涉减弱（相消）。
    \end{enumerate}
    若两波源初相相同（$\varphi_1=\varphi_2$），则相位差只由波程差 $\delta=r_2-r_1$ 决定：
    \[
    \delta=\pm k\lambda\ \text{时加强},\qquad
    \delta=\pm(2k+1)\frac{\lambda}{2}\ \text{时减弱}.
    \]

    \corepoint{（1）相干条件三条缺一不可，尤其是"相位差恒定"，两个独立的普通光源不能产生稳定干涉；（2）干涉加强不是"位移始终最大"，而是合振幅最大、能量集中；（3）当 $\varphi_1\ne\varphi_2$ 时，必须用相位差公式而不能只看波程差。}

\section{驻波}
    \subsection{驻波的形成}
    两列振幅相同、传播方向相反的相干波叠加，就形成驻波。典型情形是波在介质界面上的反射波与入射波叠加。

    \subsection{驻波的波函数}
    设两列波分别沿 $+x$、$-x$ 方向传播，波函数为
    \[
    y_1=A\cos\left[2\pi\left(\nu t-\frac{x}{\lambda}\right)\right],\qquad
    y_2=A\cos\left[2\pi\left(\nu t+\frac{x}{\lambda}\right)\right],
    \]
    叠加后
    \[
    y=y_1+y_2=2A\cos\frac{2\pi x}{\lambda}\cos2\pi\nu t.
    \]
    即
    \keyformula{y=\left(2A\cos\frac{2\pi x}{\lambda}\right)\cos\omega t}
    这就是驻波的波函数。

    驻波与行波有显著不同：
    \begin{enumerate}
        \item 各点的振幅为 $\left|2A\cos\dfrac{2\pi x}{\lambda}\right|$，不随时间变化，空间上形成稳定的振幅分布；
        \item 当 $\cos\dfrac{2\pi x}{\lambda}=0$，即 $x=(2k+1)\dfrac{\lambda}{4}$ 时，振幅为零，这些点称为波节；
        \item 当 $\left|\cos\dfrac{2\pi x}{\lambda}\right|=1$，即 $x=k\dfrac{\lambda}{2}$ 时，振幅最大（$2A$），这些点称为波腹；
        \item 相邻两波节（或两波腹）之间的距离为 $\lambda/2$，相邻波节与波腹之间的距离为 $\lambda/4$；
        \item 两波节之间的各点振动相位相同，波节两侧的各点振动相位相反；
        \item 驻波不传播能量，能量在波节与波腹之间的区域内来回转换，平均能流为零。
    \end{enumerate}

    \subsection{半波损失}
    当波从波疏介质（密度较小、波速较大）入射到波密介质（密度较大、波速较小）的分界面上反射时，反射波的相位发生 $\pi$ 的突变，相当于损失了半个波长，这种现象称为半波损失。反之，从波密介质射向波疏介质而反射时，无半波损失。

    半波损失的存在使驻波的波节出现在反射界面处（波密介质一侧），这是弦线、空气柱中形成驻波时确定波节位置的重要依据。

    \subsection{弦线振动的简正模式}
    两端固定的弦（长为 $L$）上形成驻波时，两端必为波节，故弦长必须是半波长的整数倍：
    \[
    L=n\frac{\lambda_n}{2}\quad(n=1,2,3,\cdots),
    \]
    相应的波长和频率为
    \[
    \lambda_n=\frac{2L}{n},\qquad
    \nu_n=\frac{u}{\lambda_n}=\frac{n}{2L}\sqrt{\frac{F}{\rho_l}}.
    \]
    $n=1$ 对应最低频率 $\nu_1=\dfrac{u}{2L}$，称基频（基音）；$n=2,3,\cdots$ 对应的频率称谐频（泛音），都是基频的整数倍。每一种可能的稳定振动方式称简正模式。

    一端固定、一端自由的弦（或空气柱），自由端为波腹，则
    \[
    L=(2n-1)\frac{\lambda_n}{4},\qquad
    \nu_n=\frac{(2n-1)u}{4L}.
    \]
    简正模式是乐器发声、微波谐振腔、激光谐振腔等问题的基础。

    \corepoint{（1）驻波不是波，而是特殊的振动状态，它不传播能量；（2）波节处振幅为零但并非"没有振动"，各点仍按同一频率振动，只是相位不同；（3）半波损失只发生在"波疏到波密"的反射，且只对反射波有 $\pi$ 的相位突变。}

\section{多普勒效应}
    当波源和观察者（接收器）之间有相对运动时，观察者接收到的频率与波源发出的频率不同，这种现象称为多普勒效应。

    设波源频率为 $\nu$，波在介质中的速度为 $u$，波源相对介质的速度为 $v_s$，观察者相对介质的速度为 $v_0$，规定二者相互接近时 $v_s$、$v_0$ 取正，相互远离时取负。则观察者接收到的频率为
    \keyformula{\nu'=\nu\frac{u+v_0}{u-v_s}}
    可分三种情形讨论：
    \begin{enumerate}
        \item 波源静止（$v_s=0$）、观察者运动：$\nu'=\nu\dfrac{u+v_0}{u}$。观察者接近波源时 $v_0>0$，接收频率升高；远离时 $v_0<0$，接收频率降低；
        \item 观察者静止（$v_0=0$）、波源运动：$\nu'=\nu\dfrac{u}{u-v_s}$。波源接近观察者时 $v_s>0$，波长被"压缩"，接收频率升高；远离时频率降低；
        \item 波源与观察者同时运动：由上式统一给出。
    \end{enumerate}
    注意两种情形的不对称性：波源运动时波长改变，观察者运动时波长不变而"接收速率"改变，因此同样速度下二者引起的频率变化并不相同。

    当波源速度超过波速（$v_s>u$）时，波源前方各子波波前的包迹形成圆锥面，称马赫锥，其表面就是冲击波，伴随巨大的声响（如超音速飞机的"音爆"）。

    光也有多普勒效应。由于光速与参考系无关，光的 Doppler 效应需用相对论公式处理，其结果是：光源与观察者相互远离时接收频率降低（光谱线向红端移动，称红移），相互接近时频率升高（紫移）。天文上正是通过测量星系光谱的红移，得出星系正在远离我们（哈勃定律），成为宇宙膨胀的重要证据。

    \corepoint{多普勒效应中"波源运动"与"观察者运动"的效果不同，不能简单互换；公式中的符号必须按"接近为正"统一规定，否则结果会颠倒。}

\section{声波}
    声波是机械纵波，频率在 $20\ \mathrm{Hz}\sim20\ \mathrm{kHz}$ 之间能引起人耳听觉的称可闻声波；低于 $20\ \mathrm{Hz}$ 的称次声波，高于 $20\ \mathrm{kHz}$ 的称超声波。

    声波在介质中传播时，介质各点的压强在原来压强的基础上发生周期性变化，这一压强的变化量称声压：
    \[
    p=p_m\cos\left[\omega\left(t-\frac{x}{u}\right)+\varphi\right],
    \]
    声压幅值 $p_m=\rho uA\omega$。声强（能流密度）为
    \[
    I=\frac{1}{2}\rho A^{2}\omega^{2}u=\frac{p_m^{2}}{2\rho u}.
    \]
    人耳对声波强弱的感受并不与声强成正比，而是近似与声强的对数成正比，故引入声强级：
    \[
    L=10\lg\frac{I}{I_0}\ \ (\mathrm{dB}),
    \]
    其中 $I_0=10^{-12}\ \mathrm{W/m^{2}}$ 为基准声强（约相当于人耳能听到的最弱声音）。

    超声波频率高、波长短，具有方向性好、穿透力强、能量集中的特点，可用于探伤、测距、医学成像（B 超）和清洗；次声波频率低、波长长、传播远、衰减小，可用于监测地震、核爆炸，但也可能对人体造成危害。

    \corepoint{声压与声强是两个不同的量，声压是压强的变化量（单位 $\mathrm{Pa}$），声强是能流密度（单位 $\mathrm{W/m^{2}}$）；声强级用分贝表示，是相对量。}

\section{电磁波}
    电磁波是变化的电场和变化的磁场在空间的传播，其理论由麦克斯韦方程组给出（详见第八章）。电磁波的性质与机械波有相似之处，也有本质区别：
    \begin{enumerate}
        \item 电磁波是横波，$\boldsymbol{E}\perp\boldsymbol{B}\perp\boldsymbol{u}$，三者构成右手螺旋关系，且 $\boldsymbol{E}$ 与 $\boldsymbol{B}$ 同相位；
        \item 电磁波在真空中的传播速度为 $c=\dfrac{1}{\sqrt{\varepsilon_0\mu_0}}\approx3.00\times10^{8}\ \mathrm{m/s}$，在介质中 $u=\dfrac{c}{\sqrt{\varepsilon_r\mu_r}}$；
        \item 电磁波不需要介质，可以在真空中传播；
        \item 电磁波具有能量和动量，其能流密度由坡印廷矢量 $\boldsymbol{S}=\boldsymbol{E}\times\boldsymbol{H}$ 描述；
        \item 电磁波的传播速度、波长与频率同样满足 $u=\lambda\nu$，进入介质时频率不变、波速和波长改变。
    \end{enumerate}

    电磁波谱按波长由长到短依次为：无线电波、微波、红外线、可见光、紫外线、$X$ 射线、$\gamma$ 射线。它们本质上都是电磁波，只是波长（频率）不同，因而产生方式和与物质相互作用的方式不同。可见光只占电磁波谱中极窄的一段（约 $400\sim760\ \mathrm{nm}$）。

    \corepoint{电磁波与机械波的最大区别在于：电磁波不需要介质，可在真空中传播；且电磁波是横波，其"振动"是电场和磁场的变化，而不是介质质点的位移。}
